% sec_params.tex
\section{Параметры и определения}

Для описания эфира и составляющих его частиц будут использоваться следующие величины и параметры. Частицы эфира, Амеры, — это абсолютно упругие, гладкие сферы диаметром \(d_a\) и массой \(m_a\). Длина их свободного пробега в среде Свободного эфира \(\lambda_{\text{сэ}}\) равна среднему значению модуля перемещения между парами последовательных соударений. Аналогично, для частиц в направленном потоке вихря длина свободного пробега \(\lambda_{\text{эв}}\). Концентрация частиц в Свободном эфире \(n_{\text{сэ}}\), а в потоке эфира вихря \(n_{\text{эв}}\). Соответственно скорости \(\vec{v}_{\text{сэ}}\), \(\vec{v}_{\text{эв}}\), а импульсы \(\vec{P}_{\text{сэ}} = m_a \vec{v}_{\text{сэ}}\) и \(\vec{P}_{\text{эв}} = m_a \vec{v}_{\text{эв}}\).

«Заполненность» мирового пространства эфиром можно охарактеризовать параметром \(\frac{\lambda_{\text{сэ}}}{d_a}\), позволяющим оценить совокупную значимость пространственной геометрии частиц эфира в процессе осуществления объектов весомой материи. В существующих, известных автору, моделях эфира приводятся различные значения этого параметра. Наибольшее значение параметра \(\frac{\lambda_{\text{сэ}}}{d_a} \sim 10^{30}\) приведено у В.А. Ацюковского в его «Началах эфиродинамического естествознания» (М., кн.2, 2009). Для получения представления о том, что это означает, воспользуемся аналогией. Условно предположим, что \textit{Амер} имеет диаметр один метр, тогда длина его свободного пробега будет сопоставима с диаметром Солнечной системы. Наименьшее значение величины параметра \(\frac{\lambda_{\text{сэ}}}{d_a} \sim 10^{8}\) приведено у Бычкова В.Л. и Зайцева Ф.С. (Математическое моделирование электромагнитных и гравитационных явлений, М., МГУ, 2019), где эфир представлен как аналог «\textbf{сыпучей}» среды. Использование аналогии говорит, что «метровый» \textit{Амер} между последовательными соударениями пролетает расстояние от Земли до Луны, что никак не тянет на «сыпучесть». Приведённые примеры параметра \(\frac{\lambda_{\text{сэ}}}{d_a}\) указывают на то, что модельные представления авторов предполагают малую заполненность пространства Свободным эфиром, из чего можно сделать вывод о том, что они не рассматривали вопросы осуществления вихревых объектов весомой материи. Столь значительные отличия величины длины свободного пробега частицы и её диаметра не позволяют составить физический сценарий сохранения целостности вихревого объекта. Иными словами, «в таком "решете" воду не унесёшь», хотя каждый физик знает, что всё зависит от размера ячейки. Таким образом, есть основания предполагать, что значение параметра \(\frac{\lambda_{\text{сэ}}}{d_a} \ll 10^{8}\).

Из вышеизложенного следует, что кинетические параметры и концентрация частиц Свободного эфира и вихревого потока объектов весомой материи не могут иметь случайную величину, а должны быть взаимно обусловленными или детерминантным образом связанными. Иными словами, существование объектов весомой материи возможно при условии, что скорости хаотического движения частиц Свободного эфира и их концентрация связаны со скоростями, направлением движения и концентрацией частиц вихревого потока в рамках физического сценария осуществления объектов. Исходя из этого, попытаемся сформулировать условия, соблюдение которых позволяет эфирным вихревым объектам весомой материи сохранять свою целостность.

Вихревой объект весомой материи — это замкнутая, глубоко симметричная, равновесная, эфирокинетическая система взаимодействующих посредством соударений Амеров как внутри вихревого объекта, так и на границе со Свободным эфиром.
